\documentclass[../../main.tex]{subfiles}% 注意这里的文件路径不能用 ./main.tex ,否则用latexmk编译子文件会报错
\graphicspath{{\subfix{./image/}}} % 指定图片目录,后续可以直接使用图片文件名
% 注意这里的文件路径不能用 ../../image/ ,否则用latexmk编译子文件会报错

% 例如:
% \begin{figure}[H]
% \centering
% \includegraphics[scale=0.3]{图.png}
% \caption{}
% \label{figure:图}
% \end{figure}
% 注意:上述\label{}一定要放在\caption{}之后,否则引用图片序号会只会显示??.

\begin{document}

\section{一个方程的群}

\begin{theorem}
设$f(x)$是域$F$上无重根的首一多项式，$K$为$f(x)$的分裂域，在$K$上$f(x)=\prod\limits_{i=1}^m\limits (x-\alpha_i)$，则
\begin{enumerate}[(1)]
\item $\mathrm{Gal}(K/F)$同构于对称群$S_{\{\alpha_1,\alpha_2,\dots,\alpha_m\}}$的一个子群；
\item $f(x)\in F[x]$不可约当且仅当$\mathrm{Gal}(K/F)$在$\{\alpha_1,\alpha_2,\dots,\alpha_m\}$上的作用可递。
\end{enumerate}
\end{theorem}

\begin{proof}

记$X=\{\alpha_1,\alpha_2,\dots,\alpha_m\}$。
(1) 对任意$\sigma\in\mathrm{Gal}(K/F)$，由$f(\alpha_i)=0$，得$f(\sigma(\alpha_i))=\sigma(f(\alpha_i))=0$，故$\sigma(\alpha_i)\in X$，即$\sigma|_X\in S_X$。
由$K=F(\alpha_1,\alpha_2,\dots,\alpha_m)$，$\sigma|_X=\tau|_X$当且仅当$\sigma=\tau$，因此$\mathrm{Gal}(K/F)\cong\{\sigma|_X\mid \sigma\in\mathrm{Gal}(K/F)\}\leqslant S_X$。

(2) 必要性：若群作用可递，则对任意$\alpha_i,\alpha_j\in X$，存在$\sigma\in\mathrm{Gal}(K/F)$使得$\sigma(\alpha_i)=\alpha_j$，故$\mathrm{Irr}(\alpha_i,F)=\mathrm{Irr}(\alpha_j,F)$。于是$f(x)=\prod\limits_{i=1}^m\limits (x-\alpha_i)\mid \mathrm{Irr}(\alpha_i,F)$，即$f(x)=\mathrm{Irr}(\alpha_i,F)$，故$f(x)$不可约。

充分性：若$f(x)$不可约，则$\mathrm{Irr}(\alpha_i,F)=\mathrm{Irr}(\alpha_j,F)$，存在$F$-同构$\sigma'\colon F(\alpha_i)\to F(\alpha_j)$满足$\sigma'(\alpha_i)=\alpha_j$。由同构开拓定理，$\sigma'$可开拓为$K$的$F$-自同构$\sigma\in\mathrm{Gal}(K/F)$，即$\sigma(\alpha_i)=\alpha_j$，故群作用可递.

\end{proof}

\begin{remark}
记$\mathrm{Gal}(K/F)$为多项式$f(x)$（或方程$f(x)=0$）对域$F$的群，简记为$G_f$或$G(f,F)$，也可记为$\mathrm{Gal}(K/F)$。由Cayley定理，任意有限群均可视为对称群的子群，故只需证明对$S_n$存在对应多项式，即可推出对一般有限群存在对应多项式。
\end{remark}

\begin{theorem}
设$F$为域，$x_1,x_2,\dots,x_n$为不定元，$p_1,p_2,\dots,p_n$为$x_1,x_2,\dots,x_n$的初等对称多项式，记$g(x)=\prod\limits_{i=1}^n\limits (x-x_i)\in F(p_1,p_2,\dots,p_n)[x]$，则
$$\mathrm{Gal}(F(x_1,x_2,\dots,x_n)/F(p_1,p_2,\dots,p_n))\cong S_n.$$
\end{theorem}

\begin{proof}

记$G=\mathrm{Gal}(F(x_1,x_2,\dots,x_n)/F(p_1,p_2,\dots,p_n))$。
对任意$\sigma\in S_n$，$\sigma$可延拓为多项式环$F[x_1,x_2,\dots,x_n]$的$F$-自同构：$\sigma(a)=a\ (\forall a\in F)$，$\sigma(x_i)=x_{\sigma(i)}$，进一步延拓为分式域$F(x_1,x_2,\dots,x_n)$的$F$-自同构，故$S_n\leqslant G$。
由$\sigma(p_i)=p_i\ (1\leqslant i\leqslant n)$，得$F(p_1,p_2,\dots,p_n)\subseteq \mathrm{Inv}S_n$，因此
\begin{align*}
[F(x_1,x_2,\dots,x_n):\mathrm{Inv}S_n]\leqslant [F(x_1,x_2,\dots,x_n):F(p_1,p_2,\dots,p_n)].
\end{align*}
$F(x_1,x_2,\dots,x_n)$是$g(x)$的分裂域，$\deg g(x)=n$，故$[F(x_1,x_2,\dots,x_n):F(p_1,p_2,\dots,p_n)]\leqslant n!$。
结合$[F(x_1,x_2,\dots,x_n):\mathrm{Inv}S_n]\geqslant n!$，得
\begin{align*}
[F(x_1,x_2,\dots,x_n):\mathrm{Inv}S_n]=[F(x_1,x_2,\dots,x_n):F(p_1,p_2,\dots,p_n)]=n!,
\end{align*}
于是$\mathrm{Inv}S_n=F(p_1,p_2,\dots,p_n)$，即$G\cong S_n$.

\end{proof}

\begin{corollary}
$g(x)=\prod\limits_{i=1}^n\limits (x-x_i)$是$F(p_1,p_2,\dots,p_n)$上的不可约多项式。
\end{corollary}

\begin{proof}

$S_n$在$\{x_1,x_2,\dots,x_n\}$上的作用可递，故$g(x)$不可约.

\end{proof}

\begin{example}
多项式$x^3-2\in\mathbb{Q}[x]$的分裂域为$E=\mathbb{Q}(\sqrt[3]{2},\omega)$，其中$\omega$是$x^2+x+1$的根，$[E:\mathbb{Q}]=6$。
$\mathbb{Q}(\sqrt[3]{2}),\mathbb{Q}(\omega\sqrt[3]{2}),\mathbb{Q}(\omega^2\sqrt[3]{2})$均为$E$的3次中间域，$\mathrm{Gal}(E/\mathbb{Q})$含3个2阶子群，为非Abel的6阶群，因此$\mathrm{Gal}(E/\mathbb{Q})\cong S_3$。
\end{example}

\begin{lemma}
设$p$为素数，若$S_p$的子群$G$包含一个$p$阶元与一个对换，则$G=S_p$。
\end{lemma}

\begin{proof}

由Sylow第二定理，$S_p$的Sylow $p$-子群均为$p$阶且互相共轭。
不妨设$p$阶元$\sigma=(1\ 2\ \dots\ p)\in G$，对换$(1\ 2)\in G$，则
\begin{align*}
\sigma^k(1\ 2)\sigma^{-k}=(k+1\ k+2),\quad 1\leqslant k\leqslant p-2.
\end{align*}
由对称群的生成性质，上述对换生成$S_p$，故$G=S_p$.

\end{proof}

\begin{theorem}
设$p$为素数，$f(x)\in\mathbb{Q}[x]$为$p$次不可约多项式，且恰有两个非实根，则$f(x)$对$\mathbb{Q}$的群$G_f\cong S_p$。
\end{theorem}

\begin{proof}

设$f(x)=\prod\limits_{i=1}^p\limits (x-\alpha_i)$，$\alpha_1,\alpha_2$为$f(x)$的非实根。
由$[\mathbb{Q}(\alpha_1):\mathbb{Q}]=p$，得$p\mid |G_f|$，故$G_f$含$p$阶元。
复共轭运算$\tau\in\mathrm{Gal}(\mathbb{C}/\mathbb{Q})$满足$\tau(\alpha_1)=\alpha_2,\tau(\alpha_2)=\alpha_1,\tau(\alpha_i)=\alpha_i\ (i>2)$，故$\tau$在根集上的限制为一个对换。
由引理可得$G_f\cong S_p$.

\end{proof}

\begin{example}
取充分大的素数$p$，$f(x)=x^5-px+p\in\mathbb{Q}[x]$。
由$f'(x)=5x^4-p$，$f''(x)=20x^3$，$\lim\limits_{x\to-\infty}f(x)=-\infty$，$\lim\limits_{x\to+\infty}f(x)=+\infty$，可知$f(x)$恰有3个实根、2个非实根，故$G_f=S_5$。
\end{example}

\begin{example}
设$k>3$为奇数，$m$为正偶数，$n_1<n_2<\dots<n_{k-2}$为$k-2$个偶数，令
\begin{align*}
g(x)=(x^2+m)(x-n_1)(x-n_2)\cdots(x-n_{k-2}),\quad f(x)=g(x)-2.
\end{align*}
分析可得$f(x)$至少有$k-2$个实根；当$m$充分大时，$\sum\limits_{i=1}^k\limits r_i^2<0$，故$f(x)$恰有2个非实根。
由Eisenstein判别法，$f(x)$不可约，因此$G_f=S_k$。
\end{example}

\begin{remark}
对任意有限群$G$，是否存在$f(x)\in\mathbb{Q}[x]$使得$G(f,\mathbb{Q})\cong G$的问题已提出百年，至今未完全解决。
\end{remark}



\end{document}